% !TeX root = ../papers/nc-p3.tex
\subsection{三次元接点の基底が量子多項式環となる有限四辺の分類}
\label{sec:quantum-tangent-windows}
前節の $B$ が量子多項式環
$B_Q=k\langle x_0,x_1,x_2\rangle/(x_jx_k-q_{jk}x_kx_j)$ になる枝を扱う。
$q_{jk}q_{kj}=q_{jj}=1$ とし、$q_{01}=p,q_{12}=q,q_{20}=r$ とする。
この基底は順序付き単項式の基底と量子外積双対を持つ三次元Koszul左右AS正則代数である。

\begin{theorem}\label{thm:quantum-tangent-finite}
有限四辺で三次元接点からこの形の $B_Q$ が得られる全AS対象は、
次の正規表示で尽くされる。添字 $a,b,c$ は $0,1,2$ の巡回順である。
\begin{gather*}
 p=q=r=\zeta,\qquad \zeta=-1\ \text{または}\ \zeta^2=-1,\\
 X_n x_a-x_a X_{n+1}
  =\sum_{s=0}^2 C_{a,s}\,\zeta^{n(s+1)}x_b^s x_c^{2-s}.
\end{gather*}
全 $C\in k^{3\times3}$ が標準AS正則である。四辺が有限である必要十分条件は
\[
 (C_{a,0},C_{a,1},C_{a,2})\ne0\quad(\forall a),\qquad
 (C_{j-1,2},C_{j-2,0})\ne(0,0)\quad(\forall j)
\]
であり、行添字は法3で読む。三次元接点は一意である。
\end{theorem}
\begin{proof}
量子トーラスで $x_j x^\alpha=\chi_j x^\alpha x_j$、
$\chi_j=\prod_k q_{jk}^{\alpha_k}$ と書く。
写像 $d_n$ を多次数で分け、$d_n(x_j)=c_j(n)x^\alpha x_j$ と置く。
$\alpha+e_j$ が非負でない場合は $c_j(n)=0$ とする。
式\eqref{eq:ternary-shift-derivation}は共通の量子単項式の係数を取ると
\[
          c_j(n)-c_k(n)=\chi_k c_j(n+1)-\chi_j c_k(n+1)
\]
となる。$\alpha=e_a$ のとき、$\chi=(\chi_0,\chi_1,\chi_2)$ と
$\mathbf1=(1,1,1)$ が独立なら、
$\mathbf1\wedge k^3$ と $\chi\wedge k^3$ の交叉は
$k(\mathbf1\wedge\chi)$ であり、$c(n)$ は $\langle\mathbf1,\chi\rangle$ に入る。
これは $L_H+R_H$ のこの次数の部分である。
比例する場合 $\chi_a=1$ なので $x_a$ は中心的で、
$k^3/k\mathbf1$ の二次元が残り、移動は恒等である。

他の次数は $\alpha=-e_a+s e_b+(2-s)e_c$、$s=0,1,2$ であり、
ただ一つの成分 $c_a$ が許される。
非零な両側列の存在は $\chi_b=\chi_c$ に同値で、成立すれば
$c_a(n+1)=\chi_b^{-1}c_a(n)$ となる。
条件は $q_{bc}^2=q_{ba}/q_{ca}$ で、$s$ によらない。
成立する軸では三次元が残り、この次数に補生成元の変更は作用しない。
従ってこの計算は任意の非零 $p,q,r$ に対する $P_{B_Q}$ も決めている。

軸 $a$ の条件が破れれば、$x_a=0$ 上に制限した $d_n(x_a)$ は恒等的に零となる。
実際、非負多次数の項は $x_a$ を含み、負の次数の項は排除される。
座標直線 $[u_b:u_c]$ にある基底の点とその次点を
$u=(u_b,u_c)$、$v=(u_b,q_{bc}u_c)$ とする。混合関係は
二つの追加座標 $z,z'$ について $zv-uz'=d_n(u,v)$ である。
$q_{bc}\ne1$ なら稠密開集合で二行が独立になり全点が持ち上がる。
$q_{bc}=1$ なら二行の両立条件は二元三次式一本であり、
射影直線上の少なくとも一点で解が一次元になる。
従って有限四辺には三軸全ての条件が必要である。
その三式は $pqr\cdot q=pqr\cdot r=pqr\cdot p=1$ なので、
$p=q=r=\zeta$, $\zeta^4=1$ に限る。

$\zeta=1$ なら射影平面上の $u=v$ について $d_n(u,u)$ は三つの二次式である。
共通零点があればそこで、なければこれらが定める次数2の自己射の固定点で、
$d_n(u,u)$ は $u$ に比例する。固定点の存在は射影平面の積でグラフと対角の交叉数が
$1+2+4=7$ であることから従う。そこで $z,z'$ に一次元の解があり有限でない。
残る $\zeta$ には中心的な一次生成元がなく、負の次数の三ブロックだけが残る。
$\chi_b^{-1}=\zeta^{s+1}$ となり、単項式の定数を吸収すれば定理の式を得る。
前節によって全係数のAS正則性が従う。

$\zeta^3\ne1$ のため、三つの純粋関係を満たす基底の射影零点対は
座標三角形上のグラフである。三座標が全て非零なら巡回積で
$\zeta^3=1$ を得て矛盾し、各直線では先の $u,v$ の表示になる。
スキームとしても二座標非零の開部分では、二つの式で標的を消去した残る式が
$(1-\zeta^3)u_0u_1u_2=0$ の単元倍である。
基底への射影が一方の位置で零なら混合関係から両方で零となり、指定点だけを得る。
他では軸 $a$ を欠く直線の開トーラス上で $z,z'$ が一意に解ける。残る式は始点0で
\[
 g_a(u_b,u_c)=C_{a,2}u_b^2+\zeta C_{a,1}u_bu_c+\zeta C_{a,0}u_c^2=0
\]
である。頂点 $j$ は他の二つの混合関係に現れる $x_j^2$ の係数がともに零なら
一次元の持上げを持ち、一方が非零なら持上げを持たない。
各始点で非零係数の位置は同じなので、主張した二条件を得る。
指定点で混合三式の微分は階数3、他の局所環は直線上の二次式の根の局所環なので
三次元接点は一意である。
\end{proof}

\begin{corollary}\label{cor:quantum-tangent-invariants}
有限な正規表示の添字保持同型は以下の作用で厳密に決まる。
対角基底変更 $(t_0,t_1,t_2)$ は各 $C_{a,s}$ に
$t_b^s t_c^{2-s}/t_a$ を掛ける。偶置換は三つの行を巡回させる。
奇置換は $\zeta$ を $\zeta'=\zeta^{-1}$ に送り、対応する行で
$s$ 列を $2-s$ 列に移し、係数に $(\zeta')^{-s(2-s)}$ を掛ける。
従って原始四乗根は一方だけ採用すればよく、$-1$ 型と原始四乗根型は非同型である。
各型は既約な六次元の同型類の族で、一般最小周期はそれぞれ2と4である。
\end{corollary}
\begin{proof}
一意な点から $B_Q$ と $P_{B_Q}$ が復元される。
基底の次数付き同型は点スキームの座標三角形を保つので単項行列であり、
関係を代入すると示した置換だけが可能になる。
奇置換の係数は $x_{c'}^s x_{b'}^{2-s}$ を巡回順に戻す交換係数である。
従って指定点を外しても群作用の条件は必要十分である。
任意の零係数を許して、有限個の置換を試した後、非零座標の台を比較し、
重み行列の転置の整数核から来るLaurent単項式を比較すればよい
（定理\ref{thm:curvilinear-window-orbits}のトーラスの証明）。
移動は各列への $\zeta^{s+1}$ 倍なので、この判定を
$j=1,2,4$ 回の移動に順に用いれば特殊値の最小周期も決まる。

有限条件は九次元係数空間の非空開集合である。全係数が非零なら
重みの差により三つの軸の差を張り、各重みの成分和が1なので重み階数は3である。
従って像の次元は $9-3=6$。一般には三つの比
$C_{a,1}^2/(C_{a,0}C_{a,2})$ は互いに異なる。
これらは移動と対角変換で不変で、同じ型の置換で行に沿って移る
（$\zeta=-1$ の逆順の符号も平方で消える）。従って一般点の周期判定では置換は恒等である。
隣り合う列の重み比から、$j$ 移動を実現するには全巡回順で
$t_b/t_c=\zeta^j$ が必要で、三つを掛けると $\zeta^{3j}=1$ を得る。
四乗根なので最小の $j$ はその位数であり、そこで表示そのものが戻る。
\end{proof}

有限な各始点で、$S_a=\{s:C_{a,s}\ne0\}$ と置けば指定点の長さは
\[
                  20-\sum_a(\max S_a-\min S_a)
\]
である。実際、各開トーラスでの根の重複度の和は表示した幅であり、
頂点の上に点はなく、全長は20である。
幅2の二根は $\zeta^2 C_{a,1}^2-4\zeta C_{a,0} C_{a,2}=0$ で二重点になる。
始点を移すと判別式に非零定数が掛かるだけである。
端の係数を零にして有限条件を保つ特殊化では、失われた開トーラス上の長さが
指定点に加わる。全ての行で第0列だけが非零の場合は長さ20の一点となり最小周期は1である。
原始四乗根で三行を $(2,0,5),(7,0,13),(17,0,23)$ とすれば最小周期2の例になる。
\path{computations/verify_quantum_tangent_windows.py} はこれらの同型・周期判定と
$P_B$ の独立な行列計算を検査する。九係数の512通りの台のうち有限なものを
指定点の長さ14から20の順に数えると $8,24,54,50,36,6,2$ 通りである。
この数は標識付き台の数であり、同型類数ではない。
